\paragraph{Joint deleted profiles can require branching.}
The compatible face erasures do not give a one-source reconstruction of
all selected profiles, even if arbitrary passive helpers are allowed.

\begin{proposition}[Two source copies for joint profile requirements]
\label{core:prop:joint-profile-repetition}
Fix distinct owners $a,b,c$, and put
\[
 H=((Z\star_aY)\star_bW)\star_c((Z\star_bV)\star_aW),
 \qquad C=H+Z.
\]
Among core words $S\subseteq C$ whose singleton profiles after excluding
$Z$ are
\[
 \Phi_a(S)|_{Z=0}=YW,\qquad
 \Phi_b(S)|_{Z=0}=VW,\qquad
 \Phi_r(S)|_{Z=0}=0\quad(r\notin\{a,b\}),
\]
the minimum number of $Z$-occurrences is two. This holds in the primitive
signature and with union named. The displayed primitive word $H$ attains
the minimum.
\end{proposition}
\begin{proof}
Profile recursion gives the stated deleted profiles of $H$.
A word omitting $Z$ has a positive top profile at every owner, so cannot
meet the zero profiles. Suppose instead that $S$ contains $Z$ once.
Follow its unique path, treating each sibling as an arbitrary passive
core helper. Only owners $a,b$ can occur on that path: at the top passive
coalition a gate of any other owner would install a positive deleted
profile, and every later helper has positive top profile, so that bit
could not disappear. Union gates may also occur.

The following four inputs give the bounds used below:
\[
\begin{array}{c|cccc|c}
 &Z&Y&V&W&C\\\hline
 P_a&a&b&b&a&\{a\}\\
 P_b&b&a&a&b&\{b\}\\
 Q_b&b&a&c&\{a,b\}&\{a,b\}\\
 Q_a&a&c&b&\{a,b\}&\{a,b\}
\end{array}
\]
Letters in the input columns denote singleton menus.

At the passive input of $P_a$, the spine maps $Z=\{a\}$ to $\{a\}$
but does not map $Z=I\setminus\{a\}$ to that atom, because its deleted
$a$-profile is false on $\{W\}$. No helper can therefore be a reset
to its gate's owner: that would make the entire one-$Z$ context constant.
Every step is now extensive, including union, so the spine stays at
$\{a\}$ throughout this input. Every $b$-helper must equal $\{a\}$,
while no $a$-helper equals $\{a\}$. The symmetric argument at $P_b$
gives, for every helper at the indicated gate,
\[
\begin{array}{c|cc}
 \text{gate owner}&\Phi_a(\{W\})&\Phi_b(\{W\})\\\hline
 a&0&1\\
 b&1&0
\end{array}
\]

At the passive input of $Q_b$, the atomic sets are
$E_a=\{Y\}$, $E_b=\varnothing$, $E_c=\{V\}$.
Every $a$-helper is ambiguous. Its positive $b$-coalition $\{W\}$
excludes an $a$-singleton on the disjoint coalition $\{Y\}$ and a
$c$-singleton on $\{V\}$ by common source; no other singleton can
occur. Hence no $a$-gate resets here, and no $b$-gate resets because
there is no $b$-atomic input. By the first $a$-gate the running value
is ambiguous and still contains $b$; no later step can remove a winner.

Some $b$-helper must have positive profile on $\{V,W\}$, since only
a $b$-gate can introduce the required positive deleted bit from its
initial zero value. That helper has a negative bit on $\{W\}$ by the
table. The pair test \eqref{core:eq:profile-pair-safe} forces it to
contain $c$ at $Q_b$, using provenance: the difference input $V=\{c\}$ is outside
$\{a,b\}$, there is no $b$-atom, and the allowed foreign $a$-atom
is at $Y$, outside $\{V,W\}$. It must therefore occur before the first
$a$-gate, or the spine would retain a forbidden $c$.

Symmetrically, at $Q_a$ every $b$-helper is ambiguous, and an $a$-helper
positive on $\{Y,W\}$ and negative on $\{W\}$ must contain $c$.
Such a helper exists and must occur before the first $b$-gate.
Thus each owner must occur before the first occurrence of the other,
a contradiction. Union steps are extensive throughout and cannot remove
either forbidden contribution. This excludes every one-$Z$ word in both
signatures, completing the sharp lower bound.
\end{proof}
